% from ega0-1-fr.tex
\oldpage[0\textsubscript{I}]{35}

\begin{env}[3.8.1]
\label{0.3.8.1-fr}
Soit $X$ un espace topologique dont la topologie admet une base
$\mathfrak{B}$ formée d'ensembles ouverts \emph{quasi-compacts}. Soit
$\mathcal{F}$ un \emph{faisceau d'ensembles} sur $X$~; si on munit chacun des
$\mathcal{F}(U)$ de la topologie \emph{discrète},
$U\mapsto\mathcal{F}(U)$ est un \emph{préfaisceau d'espaces topologiques}.
Nous allons voir qu'il existe un \emph{faisceau d'espaces topologiques
$\mathcal{F}'$ associé à $\mathcal{F}$} (3.5.6) tel que
$\Gamma(U,\mathcal{F}')$ soit l'espace discret $\mathcal{F}(U)$ pour tout
ouvert \emph{quasi-compact} $U$. Il suffira pour cela de montrer que le
préfaisceau $U\mapsto\mathcal{F}(U)$ d'espaces topologiques discrets
\emph{sur $\mathfrak{B}$} vérifie la condition $(\mathrm{F}_0)$ de (3.2.2),
et plus généralement que si $U$ est un ouvert quasi-compact et si
$(U_\alpha)$ est un recouvrement de $U$ par des ensembles de $\mathfrak{B}$,
la topologie la moins fine $\mathcal{T}$ sur $\Gamma(U,\mathcal{F})$ rendant
continues les applications
$\Gamma(U,\mathcal{F})\to\Gamma(U_\alpha,\mathcal{F})$ est la topologie
\emph{discrète}. Or, il existe un nombre fini d'indices $\alpha_i$ tels que
$U=\bigcup_i U_{\alpha_i}$. Soit $s\in\Gamma(U,\mathcal{F})$ et soit $s_i$ son
image dans $\Gamma(U_{\alpha_i},\mathcal{F})$~; l'intersection des images
réciproques des ensembles $\{s_i\}$ est par définition un voisinage de $s$
pour $\mathcal{T}$~; mais puisque $\mathcal{F}$ est un faisceau d'ensembles et
que les $U_{\alpha_i}$ recouvrent $U$, cette intersection se réduit à $s$,
d'où notre assertion.

On notera que si $U$ est un ouvert non quasi-compact de $X$, l'espace
topologique $\Gamma(U,\mathcal{F}')$ a encore $\Gamma(U,\mathcal{F})$ comme
ensemble sous-jacent, mais sa topologie n'est pas discrète en général~:
c'est la moins fine rendant continues les applications
$\Gamma(U,\mathcal{F})\to\Gamma(V,\mathcal{F})$, pour $V\in\mathfrak{B}$ et
$V\subset U$ (les $\Gamma(V,\mathcal{F})$ étant discrets).

Les considérations précédentes s'appliquent sans modification aux faisceaux
de groupes ou d'anneaux (non nécessairement commutatifs), et leur associent
respectivement des faisceaux de \emph{groupes topologiques} ou d'\emph{anneaux
topologiques}. Pour abréger, nous dirons que le faisceau $\mathcal{F}'$ est le
faisceau d'\emph{espaces} (resp. \emph{groupes}, \emph{anneaux})
\emph{pseudo-discrets} associé au faisceau d'ensembles (resp. groupes,
anneaux) $\mathcal{F}$.
\end{env}

\begin{env}[3.8.2]
\label{0.3.8.2-fr}
Soient $\mathcal{F}$, $\mathcal{G}$ deux faisceaux d'ensembles (resp. groupes,
anneaux) sur $X$, $u:\mathcal{F}\to\mathcal{G}$ un homomorphisme. Alors $u$
est aussi un homomorphisme \emph{continu} $\mathcal{F}'\to\mathcal{G}'$, en
désignant par $\mathcal{F}'$ et $\mathcal{G}'$ les faisceaux pseudo-discrets
associés à $\mathcal{F}$ et $\mathcal{G}$~; cela résulte en effet de
(3.2.5).
\end{env}

\begin{env}[3.8.3]
\label{0.3.8.3-fr}
Soient $\mathcal{F}$ un faisceau d'ensembles, $\mathcal{H}$ un sous-faisceau de
$\mathcal{F}$, $\mathcal{F}'$ et $\mathcal{H}'$ les faisceaux pseudo-discrets
associés à $\mathcal{F}$ et $\mathcal{H}$ respectivement. Alors, pour tout
ouvert $U\subset X$, $\Gamma(U,\mathcal{H}')$ est \emph{fermé} dans
$\Gamma(U,\mathcal{F}')$~: en effet, c'est l'intersection des images
réciproques des $\Gamma(V,\mathcal{H})$ (pour $V\in\mathfrak{B}$,
$V\subset U$) par les applications continues
$\Gamma(U,\mathcal{F})\to\Gamma(V,\mathcal{F})$, et
$\Gamma(V,\mathcal{H})$ est fermé dans l'espace discret
$\Gamma(V,\mathcal{F})$.
\end{env}

\section{Espaces annelés}
\label{section:0.4-fr}

\subsection{Espaces annelés, $\mathcal{A}$-Modules, $\mathcal{A}$-Algèbres}
\label{subsection:0.4.1-fr}

\begin{env}[4.1.1]
\label{0.4.1.1-fr}
Un \emph{espace annelé} (resp. \emph{topologiquement annelé}) est un couple
$(X,\mathcal{A})$ formé d'un espace topologique $X$ et d'un faisceau d'anneaux
(non nécessairement commutatifs) (resp. d'un faisceau d'anneaux topologiques)
$\mathcal{A}$ sur $X$~; on dit que $X$ est l'espace
